\documentclass[11pt]{article}
\usepackage[a4paper,margin=1in]{geometry}
\usepackage{amsmath,amssymb,amsthm}
\usepackage{array}
\usepackage{hyperref}

\newcommand{\dd}{\mathrm{d}}
\newcommand{\cD}{\mathcal{D}}
\newcommand{\avg}[1]{\langle #1\rangle}
\newcommand{\Tpar}{T_{\parallel1}}
\newcommand{\Tperp}{T_{\perp1}}

\newtheorem{theorem}{Theorem}
\newtheorem{lemma}[theorem]{Lemma}
\newtheorem{proposition}[theorem]{Proposition}
\theoremstyle{remark}
\newtheorem{remark}[theorem]{Remark}

\title{Positivity of the heat-flux closure in three velocity dimensions:\\
what is proved and what is not}
\author{Working note}
\date{17 November 2026}

\begin{document}
\sloppy
\maketitle

\noindent\emph{Summary.} The goal was to prove that the quasi-static
conductivity $\chi(k)$ of the Landau operator is positive at every
wavenumber. \textbf{That is not proved here.} What is proved:
\begin{itemize}
  \item \textbf{Theorem~\ref{thm:range}.} For any conserving operator,
    $\chi(k)>0$ whenever $k^2\kappa_{\mathrm F}\eta<3$. For the Landau
    operator this is $k\lambda<0.89$, against $0.30$ before, and the same
    argument gives a cleaner upper bound on $\chi$.
  \item \textbf{Proposition~\ref{prop:red}.} An exact reduction: positivity
    is equivalent to the perpendicular temperature response not overshooting,
    $\Tperp<n_1/2$.
  \item \textbf{Theorem~\ref{thm:models}.} For the BGK operator with constant
    $\nu$ and for the Dougherty operator, in three velocity dimensions,
    $\Tpar\le\Tperp\le0$ at every $k$, hence positivity at every $k$.
\end{itemize}
For the Landau operator, numerical values satisfy $\Tpar\le\Tperp\le0$ at
every wavenumber computed, with a margin, but the computation does not reach
the collisionless regime and is not a proof. Section~\ref{sec:why} explains
why the argument that works in one velocity dimension does not extend, and
what would suffice.

\section{Reduction}

Notation is that of \texttt{landau\_3d\_closure.tex}: $g=g_e+iw$ is the
stationary solution, $n_1=\avg{g_e}>0$, $T_1=\tfrac13\avg{(|v|^2-3)g_e}$,
$\chi=-1/(k^2T_1)$, $\vartheta=v_\perp^2-2v_x^2$, $X=\avg{\vartheta g_e}$,
$\kappa_{\mathrm F}$ the conductivity and $\eta$ the viscosity. Define the
parallel and perpendicular temperature responses
\[
  \Tpar=\avg{(v_x^2-1)\,g_e},\qquad
  \Tperp=\avg{(\tfrac12v_\perp^2-1)\,g_e}.
\]

\begin{proposition}\label{prop:red}
For any operator satisfying (A3) and (S):
\begin{enumerate}
  \item[(i)] $\Tpar=-n_1<0$, $T_1=\tfrac13(\Tpar+2\Tperp)$, and
    $X=2(n_1+\Tperp)$.
  \item[(ii)] $\chi(k)>0$ if and only if $\Tperp<\tfrac12n_1$.
  \item[(iii)] If $\Tpar\le\Tperp\le0$, then
    $\tfrac13n_1\le-T_1\le n_1$, so that
    \[
      \frac{1}{k^2n_1}\;\le\;\chi(k)\;\le\;\frac{3}{k^2n_1}
      \;\le\;3\,\kappa_{\mathrm F}.
    \]
\end{enumerate}
\end{proposition}

\begin{proof}
(i) Momentum conservation gives $\avg{v_x^2g_e}=0$, hence $\Tpar=-n_1$ and
$X=\avg{v_\perp^2g_e}=2\Tperp+2n_1$. The formula for $T_1$ is the
decomposition $|v|^2-3=(v_x^2-1)+(v_\perp^2-2)$. (ii) follows from
$T_1=\tfrac13(-n_1+2\Tperp)$. (iii) uses in addition
$n_1\ge1/(k^2\kappa_{\mathrm F})$.
\end{proof}

So the sign question is whether collisions, acting on a parallel
compression, can drive the perpendicular temperature in the opposite
direction by more than half. Streaming does not act on $v_\perp$, so $\Tperp$
is produced by collisions alone.

\section{An improved range for any operator}

\begin{theorem}\label{thm:range}
Assume (A3) and (S). Then for every $k\ne0$
\begin{equation}
  -T_1\;\ge\;\frac{1}{k^2\kappa_{\mathrm F}}-\frac{\eta}{3}.
  \label{eq:lower}
\end{equation}
Consequently, if $k^2\kappa_{\mathrm F}\,\eta<3$, then $\chi(k)>0$ and
\begin{equation}
  \chi(k)\;\le\;\frac{\kappa_{\mathrm F}}{1-\tfrac13\kappa_{\mathrm F}\eta\,k^2}.
\end{equation}
\end{theorem}

\begin{proof}
From \texttt{landau\_3d\_closure.tex}: $-T_1=n_1-\tfrac13X$;
$n_1=\cD[g_e]+\cD[w]$; $\cD[w]\ge1/(k^2\kappa_{\mathrm F})$; and
$X^2\le12\,\eta\,\cD[g_e]$. Writing $D=\cD[g_e]\ge0$,
\[
  -T_1\ \ge\ \frac{1}{k^2\kappa_{\mathrm F}}+D-\frac{2}{\sqrt3}\sqrt{\eta D}
  \ \ge\ \frac{1}{k^2\kappa_{\mathrm F}}-\frac{\eta}{3},
\]
the minimum over $D$ being attained at $D=\eta/3$.
\end{proof}

\begin{remark}
The earlier bound used $D\le C_*$ and then the worst case of the cross term;
keeping $D$ on both sides removes $C_*$ from the condition. For the Landau
operator, with the computed $\kappa_{\mathrm F}=3.9502\,\tau$ and
$\eta=0.9653\,\tau$, the condition is $k\lambda<0.887$. For the Dougherty
operator it is $k\lambda<2.68$. The constants are those of the Galerkin
computation; the theorem itself is in terms of $\kappa_{\mathrm F}\eta$.
\end{remark}

\section{Two operators for which positivity holds at every wavenumber}

\begin{theorem}\label{thm:models}
In three velocity dimensions, for the BGK operator
$\Lambda=\nu(1-\Pi)$ with constant $\nu$ and for the Dougherty operator,
\[
  \Tperp=c(k)\,T_1\quad\text{with}\quad0<c(k)<1,\qquad
  T_1=-\frac{n_1}{3-2c(k)}<0 ,
\]
so $\Tpar\le\Tperp\le0$ and $\chi(k)>0$ for every $k\ne0$. Here
$c=\avg{\nu^2/(\nu^2+k^2v_x^2)}$ for BGK and
$c=2\int_0^1(1-s)E_a(s)\,\dd s$ for Dougherty.
\end{theorem}

\begin{proof}
Let $\xi=\tfrac12v_\perp^2-1$, so $|v|^2-3=(v_x^2-1)+2\xi$, $\avg{\xi}=0$,
$\avg{\xi^2}=1$, and $\Tperp=\avg{\xi g}$.

\emph{BGK.} $\Pi g=\avg{g}+\avg{v_xg}v_x+\tfrac12T_1(|v|^2-3)$, so
\[
  g=\frac{1+\nu\,\Pi g}{\nu+ikv_x}
   =\frac{1+\nu\bigl(\avg{g}+\avg{v_xg}v_x+\tfrac12T_1(v_x^2-1)\bigr)}{\nu+ikv_x}
    +\frac{\nu\,T_1\,\xi}{\nu+ikv_x}.
\]
Only the last term has a component along $\xi$, and
$\avg{\nu/(\nu+ikv_x)}=\avg{\nu^2/(\nu^2+k^2v_x^2)}=c\in(0,1)$, so
$\Tperp=c\,T_1$.

\emph{Dougherty.} In the proof of the closed form, $\Tperp=\epsilon_1
=2T_1S_1$ with $S_1=\int_0^1(1-s)E_a<\tfrac12$.

In both cases Proposition~\ref{prop:red}(i) gives
$3T_1=-n_1+2c\,T_1$, hence $T_1=-n_1/(3-2c)$, which lies in
$(-n_1,-\tfrac13n_1)$; and $\Tperp=cT_1\in(-n_1,0)$.
\end{proof}

\begin{remark}
The mechanism is that the perpendicular-energy mode $\xi$ evolves in a closed
way: the equation for the $\xi$-component is a scalar resolvent equation
driven by $T_1$, and its solution has mean in $(0,1)$ times the drive. For
the Landau operator the function $\xi$, minus its conserved part, is not an
eigenfunction and the $\xi$-sector is not invariant, so this closure does not
occur.
\end{remark}

\section{Numerical values for the Landau operator}
\label{sec:num}

From \texttt{shared/scripts/landau\_closure/run\_landau\_closure.py},
section 5 (polynomial Galerkin, degree 16; Dougherty at degree 80).
\begin{center}
\renewcommand{\arraystretch}{1.15}
\begin{tabular}{r|rrr|rrr}
  & \multicolumn{3}{c|}{Landau} & \multicolumn{3}{c}{Dougherty} \\
  $k\lambda$ & $\Tperp/n_1$ & $-T_1/n_1$ & $N_\vartheta/(9n_1)$
     & $\Tperp/n_1$ & $-T_1/n_1$ & $N_\vartheta/(9n_1)$ \\ \hline
  0.1 & $-0.981$ & 0.988 & 0.043 & $-0.995$ & 0.997 & 0.006 \\
  0.3 & $-0.856$ & 0.904 & 0.224 & $-0.958$ & 0.972 & 0.047 \\
  0.5 & $-0.712$ & 0.808 & 0.386 & $-0.893$ & 0.929 & 0.117 \\
  1.0 & $-0.459$ & 0.639 & 0.620 & $-0.696$ & 0.798 & 0.317 \\
  2.0 & $-0.245$ & 0.497 & 0.778 & $-0.423$ & 0.615 & 0.568 \\
  3.0 & $-0.164$ & 0.443 & 0.824 & $-0.290$ & 0.527 & 0.678 \\
  5.0 & $-0.094$ & 0.396 & 0.863 & $-0.174$ & 0.450 & 0.768 \\
\end{tabular}
\end{center}
Here $N_\vartheta=\avg{\vartheta,(ikv_x+\Lambda)^{-1}\vartheta}$.
\begin{itemize}
  \item At every point $-1\le\Tperp/n_1\le0$, and $-T_1/n_1$ decreases from
    $1$ towards $\tfrac13$, the value when the perpendicular temperature does
    not respond at all.
  \item The Landau values are converged in the truncation degree to three
    digits for $k\lambda\le1$ and to about two for $k\lambda\ge2$
    (degrees 10, 12, 16 give $X/(3n_1)=0.581$, $0.601$, $0.604$ at
    $k\lambda=5$). A polynomial basis does not reach larger $k\lambda$.
  \item Inequality \eqref{eq:lower} holds at every point.
\end{itemize}

\section{Why this is not a proof, and what would suffice}
\label{sec:why}

\begin{enumerate}
  \item \textbf{The structure available.} The even part of the resolvent
    defines a non-negative quadratic form
    $N(\sigma)=\avg{\sigma,(ikv_x+\Lambda)^{-1}\sigma}$ on even sources, with
    $v_x^2$ in its kernel. Modulo that kernel there are two relevant sources,
    $1$ and $\xi$, and
    \[
      n_1=N(1),\qquad \Tperp=\text{the mixed term of }1\text{ and }\xi .
    \]
    Non-negativity bounds the mixed term by $\sqrt{N(1)N(\xi)}$ in absolute
    value and says nothing about its sign. In one velocity dimension there is
    no $\xi$, $T_1=-n_1$ is a diagonal term, and positivity is automatic.
  \item \textbf{A sufficient condition.} By the Cauchy--Schwarz inequality,
    positivity holds if $N_\vartheta<9\,n_1$. In the collisionless limit the
    ratio $N_\vartheta/(9n_1)$ tends formally to $\tfrac89$, so the margin
    is $\tfrac19$; at small $k$ it tends to $0$. The table shows it
    increasing between the two. Proving $N_\vartheta<9n_1$ uniformly would
    need an upper bound on $N_\vartheta$ and a lower bound on $n_1$ at each
    $k$, both of which have variational principles; for the Landau operator
    the trial functions require inverting $\Lambda$, which is not available
    in closed form, so this would be a computer-assisted proof with
    certified bounds.
  \item \textbf{A stronger conjecture.} The data suggest
    $\Tpar\le\Tperp\le0$ for the Landau operator at every $k$: collisions
    move the perpendicular temperature towards the parallel one and never
    past zero. This is proved above for BGK and Dougherty. I do not have an
    argument for a general operator, and I do not know whether it is true
    for every conserving, isotropic, non-negative operator.
  \item \textbf{Large wavenumber.} Any proof at all $k$ also needs the
    collisionless regime, where neither the theory in these notes nor the
    polynomial computation applies to the Landau operator.
\end{enumerate}

\section{Status}

\begin{center}
\renewcommand{\arraystretch}{1.3}
\begin{tabular}{p{6.4cm}p{8.2cm}}
  Statement & Status \\ \hline
  $\chi>0$ for $k^2\kappa_{\mathrm F}\eta<3$, any operator & proved \\
  $\chi>0$ for all $k$: BGK (constant $\nu$), Dougherty & proved \\
  $\chi>0$ for all $k$: Landau & not proved; holds numerically for
     $k\lambda\le5$; proved for $k\lambda<0.89$ given the computed
     $\kappa_{\mathrm F}$, $\eta$ \\
  $\Tpar\le\Tperp\le0$: Landau & conjecture, supported numerically \\
\end{tabular}
\end{center}
The value $0.89$ uses Galerkin values of $\kappa_{\mathrm F}$ and $\eta$,
which are lower bounds for the true values by the variational
characterization of each; a certified upper bound on
$\kappa_{\mathrm F}\eta$ would be needed to make the number rigorous. No
second reader has checked the proofs.

\end{document}
